% !TEX root = ../main.tex
\chapter{Configuració de modes i presets}\label{app:configuration}\label{app:config}
\fancyhead[RE]{\slshape \textsf{Apèndix \thechapter.\ Configuració}}
\fancyhead[LO]{\slshape \textsf{\nouppercase \rightmark}}

Aquest apèndix és una fitxa de configuració. Serveix per consultar els valors
exactes que connecten la memòria amb el repositori públic: escala de Q, modes de
qualitat, presets automàtics, llindars, etiquetes de cobertura ROI i variants
auxiliars del codi. No és una secció de resultats ni una repetició del disseny:
el Capítol~\ref{ch:control} explica per què es defineixen aquests modes, el
Capítol~\ref{ch:metodologia} explica com s'avaluen i el
Capítol~\ref{ch:resultats} presenta les mesures obtingudes.

Les comandes de reproducció no es repeteixen aquí: són a
l'Apèndix~\ref{app:reproducibility}. Aquest apèndix conté només els paràmetres
que aquelles comandes consumeixen o generen. El registre de procedència de les
xifres queda reservat a l'Apèndix~\ref{app:evidence}.

\section{Escala Q/lambda}

La configuració principal utilitza l'escala \code{lambda005}, on el valor Q de
8 bits es converteix a lambda mitjançant:

\begin{equation}
\lambda(Q)=\frac{Q}{255}\cdot0,05.
\label{eq:app_q_lambda}
\end{equation}

\begin{table}[H]
\centering
\small
\begin{tabularx}{\textwidth}{r r X}
\toprule
Q & $\lambda$ & Lectura dins la configuració \\
\midrule
64 & 0,01255 & Valor de fons molt agressiu, reservat a proves de Q baixa. \\
96 & 0,01882 & Valor de fons agressiu, també fora de la configuració principal. \\
128 & 0,02510 & Mínim utilitzat per al fons en la ruta principal. \\
204 & 0,04000 & Referència uniforme històrica i baseline principal. \\
240 & 0,04706 & Valor alt per preservar ROI en configuracions de dos nivells. \\
255 & 0,05000 & Valor màxim representable en l'escala actual. \\
\bottomrule
\end{tabularx}
\caption{Valors representatius de Q i lambda en l'escala \code{lambda005}.}
\label{tab:q_lambda_appendix}
\end{table}

El format \code{uint8} permet Q entre 0 i 255. En aquesta memòria, la ruta
principal utilitza Q128--Q255 perquè és la zona amb més cobertura de validació.
Els valors Q$<128$ no són invàlids: simplement impliquen una degradació més
agressiva i es tracten com una línia experimental que requereix validació
específica.

\section{Modes principals}

Els modes de la Taula~\ref{tab:app_main_modes} són els noms que apareixen a la
memòria i a les rutes de reproducció. La taula fixa la regla de Q; la
interpretació experimental es dona al cos principal.

\begingroup
\scriptsize
\setlength{\tabcolsep}{2pt}
\begin{longtable}{p{2.3cm}p{2.5cm}p{5.4cm}p{2.1cm}}
\caption{Modes principals de Q-map documentats a la memòria.}\label{tab:app_main_modes}\\
\toprule
Mode & Entrada de decisió & Regla de Q & Ús de reproducció \\
\midrule
\endfirsthead
\toprule
Mode & Entrada de decisió & Regla de Q & Ús de reproducció \\
\midrule
\endhead
\code{q204} & Q explícita & Tots els blocs reben Q=204. No usa ROI ni target.
& Q-map base i comparació uniforme. \\
\code{global\_target} & MSE o PSNR objectiu & La calibració calcula un únic
$Q^\star$ i tots els blocs reben aquest valor. & Ruta mesurada global. \\
\code{adaptive\_s8} & Error relatiu per bloc & Cada bloc rep
$Q_i=\operatorname{clip}(204+8z_i)$, amb mitjana propera a 204. & Ruta
mesurada adaptativa. \\
\code{focus\_bgq128} & Preset + base adaptativa & ROI:
$\operatorname{clip}(Q_{\mathrm{adaptive},i}+16)$; fons: Q128. & Ruta amb
preset i fons Q128. \\
\code{preserve-roi} & Preset + Q fixa ROI/fons & ROI: Q240 o Q255; fons: Q128.
No usa base adaptativa dins la ROI. & Rutes preserve Q240/Q255. \\
\bottomrule
\end{longtable}
\endgroup

\section{Modes auxiliars}

Hi ha variants implementades que ajuden a reproduir campanyes de cribratge, però
no són les rutes principals de lectura de la memòria:

\begin{table}[H]
\centering
\small
\begin{tabularx}{\textwidth}{p{3.0cm}X}
\toprule
Variant & Ús de reproducció \\
\midrule
\code{boost-only} & Manté la base adaptativa al fons i incrementa Q només a la
ROI. Variant disponible al generador semàntic. \\
\code{focus\_bgpen24} & Configuració auxiliar del full C-only: reforça ROI i
redueix el fons respecte de la base adaptativa, sense imposar Q128. \\
\code{target-focus} & Variant que calcula Q dins la ROI a partir d'un objectiu
de MSE o PSNR; correspon a la ruta \code{target\_focus\_bgq128\_measured} quan
el fons es fixa a Q128. \\
\bottomrule
\end{tabularx}
\caption{Modes auxiliars documentats per traçabilitat. No substitueixen els
modes principals de la Taula~\ref{tab:app_main_modes}.}
\label{tab:app_aux_modes}
\end{table}

Les tres variants comparteixen notació. Sigui $m_i$ la màscara generada pel
preset, amb $m_i=1$ dins la ROI i $m_i=0$ al fons. Sigui
$Q_{\mathrm{adaptive},i}$ la base adaptativa del bloc. La funció
$\operatorname{clip}(\cdot)$ limita el resultat al rang de Q permès.

\paragraph{\code{boost-only}.}
Aquesta variant només modifica la ROI. El fons queda igual que a la base
adaptativa:

\begin{equation}
Q_i^{boost}=
\begin{cases}
\operatorname{clip}(Q_{\mathrm{adaptive},i}+b), & m_i=1,\\
Q_{\mathrm{adaptive},i}, & m_i=0.
\end{cases}
\end{equation}

El paràmetre $b$ és el reforç de primer pla. Aquesta política és útil per
executar una variant en què només canvien els blocs marcats com a ROI.

\paragraph{\code{focus\_bgpen24}.}
Aquesta configuració combina reforç de ROI i penalització relativa del fons:

\begin{equation}
Q_i^{bgpen24}=
\begin{cases}
\operatorname{clip}(Q_{\mathrm{adaptive},i}+16), & m_i=1,\\
\operatorname{clip}(Q_{\mathrm{adaptive},i}-24), & m_i=0.
\end{cases}
\end{equation}

A diferència de \code{focus\_bgq128}, el fons no queda fixat a Q128. Cada bloc
de fons conserva la seva variació adaptativa, però amb una penalització de 24
unitats de Q.

\paragraph{\code{target-focus}.}
Aquesta variant substitueix el reforç fix de la ROI per una decisió basada en
un objectiu de qualitat. Dins la ROI, el valor $Q_{\mathrm{target}}$ es calcula
a partir d'un MSE o PSNR objectiu mitjançant la calibració fixed-quality. El
fons pot quedar fixat o penalitzat segons la configuració concreta:

\begin{equation}
Q_i^{target}=
\begin{cases}
Q_{\mathrm{target}}, & m_i=1,\\
Q_{\mathrm{fons}}, & m_i=0\ \text{si es fixa un fons},\\
\operatorname{clip}(Q_{\mathrm{adaptive},i}-p), & m_i=0\ \text{si es penalitza}.
\end{cases}
\end{equation}

La utilitat d'aquest mode és provar una ROI guiada directament per un objectiu
de qualitat, en lloc d'un reforç constant.

\section{Presets automàtics}

Els presets transformen una magnitud derivada de la imatge en una màscara ROI.
La Taula~\ref{tab:app_presets} és una taula de consulta de la configuració
pública. Els valors coincideixen amb \code{config/auto\_thresholds\_lambda005.tsv}
quan el preset forma part de la configuració ajustada. El preset \code{clouds}
també apareix al full C-only com a complement de \code{cloud\_avoid}; es
documenta aquí per coherència amb el disseny factorial.

\begingroup
\scriptsize
\setlength{\tabcolsep}{2pt}
\begin{longtable}{p{2.6cm}p{3.1cm}p{1.7cm}p{6.5cm}}
\caption{Presets automàtics i llindars de la configuració \code{lambda005}.}
\label{tab:app_presets}\\
\toprule
Preset & Magnitud i bandes & Condició & Lectura i limitació principal \\
\midrule
\endfirsthead
\toprule
Preset & Magnitud i bandes & Condició & Lectura i limitació principal \\
\midrule
\endhead
\code{cloud\_avoid} & Fracció CBY baixa, B02/B03/B04 & $\leq0,50$ &
Família núvol/no-núvol. Sense B11/SWIR és un fallback visible, no una detecció
completa de núvols. \\
\code{clouds} & Fracció CBY alta, B02/B03/B04 & $\geq0,50$ & Complement de
\code{cloud\_avoid}; comparteix la limitació de no disposar de B11/SWIR. \\
\code{local\_contrast} & Desviació visible, B02/B03/B04 & $\geq0,035$ &
Textura visible. No representa una classe física; selecciona variació local. \\
\code{high\_ndvi} & NDVI, B08/B04 & $\geq0,50$ & Vegetació amb criteri més
restrictiu que \code{vegetation}. \\
\code{vegetation} & NDVI, B08/B04 & $\geq0,40$ & Vegetació amb llindar empíric
del dataset. \\
\code{vegetation\_green} & GNDVI, B08/B03 & $\geq0,50$ & Vegetació amb banda
verda; la lectura no és idèntica a NDVI. \\
\code{chlorophyll} & NDCI, B05/B04 & $\geq0,10$ & Proxy de clorofil·la; la
interpretació física és limitada sense més context biofísic. \\
\code{low\_ndvi} & NDVI, B08/B04 & $\leq0,15$ & NDVI baix; pot agrupar sòl,
cobertes no vegetades, ombres, aigua o valors negatius. \\
\code{dark\_regions} & Mitjana visible, B02/B03/B04 & $\leq0,26$ & Brillantor
visible baixa; pot generar ROI àmplies si l'escena és globalment fosca. \\
\code{water\_body} & NDWI, B03/B08 & $\geq0,10$ & Candidat d'aigua simple, no
una classificació hidrològica completa. \\
\bottomrule
\end{longtable}
\endgroup

Els índexs normalitzats utilitzen la forma:

\begin{equation}
I(A,B)=\frac{A-B}{A+B+\epsilon},
\end{equation}

on $\epsilon$ evita divisions per zero. Les correspondències són: NDVI amb B08 i
B04; GNDVI amb B08 i B03; NDCI amb B05 i B04; NDWI amb B03 i B08. Els llindars
són paràmetres de configuració del projecte, no constants universals de
teledetecció \cite{rouse1974ndvi,gitelson1996gndvi,mishra2012ndci,mcfeeters1996ndwi}.

\section{Cobertura ROI}

La cobertura ROI és el percentatge de blocs marcats pel preset:

\begin{equation}
C=100\frac{\sum_i m_i}{1024}.
\end{equation}

\begin{table}[H]
\centering
\small
\begin{tabularx}{\textwidth}{p{2.5cm}p{3.2cm}X}
\toprule
Etiqueta & Rang de cobertura & Ús interpretatiu \\
\midrule
\code{no\_roi} & $C=0$ & Cap bloc rep tractament de ROI. \\
\code{low\_roi} & $0<C<10\%$ & ROI petita; pot ser sensible al veïnat de blocs de fons. \\
\code{mid\_roi} & $10\%\leq C\leq70\%$ & Cobertura intermèdia, útil per comparar ROI i fons. \\
\code{high\_roi} & $70\%<C<100\%$ & ROI ampla; queda poc fons per assumir degradació. \\
\code{all\_roi} & $C=100\%$ & Tota la imatge és ROI i desapareix la redistribució ROI/fons. \\
\bottomrule
\end{tabularx}
\caption{Etiquetes de cobertura utilitzades en l'anàlisi. No són presets i no
modifiquen el Q-map.}
\label{tab:app_roi_coverage}
\end{table}

Aquestes etiquetes només classifiquen casos després de calcular la màscara. No
canvien cap llindar ni cap valor Q. Són útils per detectar situacions en què una
política espacial pot perdre sentit: ROI buida, ROI total o ROI massa dispersa.

\section{Correspondència amb fitxers públics}

\begin{table}[H]
\centering
\small
\begin{tabularx}{\textwidth}{p{3.0cm}p{4.1cm}X}
\toprule
Element & Fitxer públic & Relació amb la reproducció \\
\midrule
Llindars de presets & \path{config/auto_thresholds_lambda005.tsv} &
Llegeix els llindars que fan servir les rutes amb presets automàtics. \\
Calibració Q/lambda & \path{config/fq_calibration_lambda005.tsv} &
Proporciona l'escala i coeficients necessaris per convertir objectius de
qualitat en Q. \\
Dataset RAW & \path{data/Sentinel2A_crop_test/} & Proporciona els retalls
multiespectrals per executar exemples públics. \\
Guia de reproducció & \path{docs/informe_final/AppendixA/text_app_a.tex} &
Conté les comandes; aquest apèndix només n'especifica els paràmetres. \\
\bottomrule
\end{tabularx}
\caption{Fitxers públics que materialitzen la configuració descrita en aquest
apèndix.}
\label{tab:app_public_config_files}
\end{table}

Aquests fitxers permeten reproduir la ruta pública descrita a
l'Apèndix~\ref{app:reproducibility}. Si es canvien pesos, bandes, normalització
o escala de lambda, aquests fitxers també s'han de revisar perquè el Q-map
continuï representant la qualitat prevista.
